\SourceLeaf{449}%
\par\indent Differenziert man erneut nach $s$ und beachtet, dass
\[
\frac{\partial\partial\xi}{\partial s\partial\theta}=\frac{\partial(pX)}{\partial\theta}\qquad\text{usw.}
\]
und dass
\[
X\xi'+Y\eta'+Z\zeta'=0,
\]
so ergibt sich
\[
\begin{aligned}
\frac{\partial\partial m}{\partial s^{2}}={}&p\left(\xi'\frac{\partial X}{\partial\theta}+\eta'\frac{\partial Y}{\partial\theta}+\zeta'\frac{\partial Z}{\partial\theta}\right)+p'\left(X\frac{\partial\xi}{\partial\theta}+Y\frac{\partial\eta}{\partial\theta}+Z\frac{\partial\zeta}{\partial\theta}\right)\\
={}&p\left(\xi'\frac{\partial X}{\partial\theta}+\eta'\frac{\partial Y}{\partial\theta}+\zeta'\frac{\partial Z}{\partial\theta}\right)+mp'p'\\
={}&-\left(\xi\frac{\partial X}{\partial s}+\eta\frac{\partial Y}{\partial s}+\zeta\frac{\partial Z}{\partial s}\right)\left(\xi'\frac{\partial X}{\partial\theta}+\eta'\frac{\partial Y}{\partial\theta}+\zeta'\frac{\partial Z}{\partial\theta}\right)\\
&+\left(\xi'\frac{\partial X}{\partial s}+\eta'\frac{\partial Y}{\partial s}+\zeta'\frac{\partial Z}{\partial s}\right)\left(\xi\frac{\partial X}{\partial\theta}+\eta\frac{\partial Y}{\partial\theta}+\zeta\frac{\partial Z}{\partial\theta}\right)\\
={}&\left(\frac{\partial Y}{\partial\theta}\frac{\partial Z}{\partial s}-\frac{\partial Y}{\partial s}\frac{\partial Z}{\partial\theta}\right)X+\left(\frac{\partial Z}{\partial\theta}\frac{\partial X}{\partial s}-\frac{\partial Z}{\partial s}\frac{\partial X}{\partial\theta}\right)Y\\
&+\left(\frac{\partial X}{\partial\theta}\frac{\partial Y}{\partial s}-\frac{\partial X}{\partial s}\frac{\partial Y}{\partial\theta}\right)Z.
\end{aligned}
\]

[Wird aber das zum Punkt $x,y,z$ gehörige Flächenelement
\[
mdsd\theta
\]
durch parallele Normalen auf die Hilfskugel mit dem Radius $1$ abgebildet, so entspricht ihm dort ein Flächenstück von der Größe
\[
\begin{aligned}
\Biggl\{&X\left(\frac{\partial Y}{\partial s}\frac{\partial Z}{\partial\theta}-\frac{\partial Y}{\partial\theta}\frac{\partial Z}{\partial s}\right)+Y\left(\frac{\partial Z}{\partial s}\frac{\partial X}{\partial\theta}-\frac{\partial Z}{\partial\theta}\frac{\partial X}{\partial s}\right)\\
&+Z\left(\frac{\partial X}{\partial s}\frac{\partial Y}{\partial\theta}-\frac{\partial X}{\partial\theta}\frac{\partial Y}{\partial s}\right)\Biggr\}dsd\theta;
\end{aligned}
\]
somit ist das Krümmungsmaß der Fläche im betrachteten Punkt gleich
\[
-\frac1m\frac{\partial\partial m}{\partial s^{2}}.\biggr]
\]
